{
    \small
    \renewcommand{\arraystretch}{1.3}
    \begin{longtable}{|>{\raggedright\arraybackslash}p{1.8cm}|>{\raggedright\arraybackslash}p{2.3cm}|>{\raggedright\arraybackslash}p{6.6cm}|>{\raggedright\arraybackslash}p{3.1cm}|}
    \caption{Requisiti di vincolo}\label{tab:req-vincolo}\\
    \hline
    \rowcolor[gray]{0.88} \textbf{Codice} & \textbf{Importanza} & \textbf{Descrizione} & \textbf{Fonti} \\ \hline
    \endfirsthead
    \hline
    \rowcolor[gray]{0.88} \textbf{Codice} & \textbf{Importanza} & \textbf{Descrizione} & \textbf{Fonti} \\ \hline
    \endhead
    \mbox{R1-V-O} & Obbligatorio & Il motore di elaborazione deve essere realizzato in Java con il framework Spring Boot. & \mbox{Capitolato}, \mbox{VE-1.11} \\ \hline
    \mbox{R2-V-O} & Obbligatorio & L'interfaccia amministrativa deve essere una soluzione web-based realizzata con Angular e TypeScript; l'uso di tecnologie alternative, come Next.js, non è ammesso. & \mbox{Capitolato}, \mbox{VE-1.11}, \mbox{LG-4.2} \\ \hline
    \mbox{R3-V-O} & Obbligatorio & La persistenza interna dei dati deve avvalersi di un database relazionale. & \mbox{Capitolato} \\ \hline
    \mbox{R4-V-D} & Desiderabile & Il DBMS dovrebbe essere scelto tra quelli consigliati dall'azienda proponente: SQLite, MariaDB, PostgreSQL. & \mbox{Capitolato} \\ \hline
    \mbox{R5-V-O} & Obbligatorio & L'interazione con i modelli LLM deve avvenire tramite un'interfaccia basata su standard Open API, compatibile con le API di OpenAI. & \mbox{VE-1.8}, \mbox{LG-2.2} \\ \hline
    \mbox{R6-V-O} & Obbligatorio & Il sistema deve utilizzare modelli a pesi aperti eseguibili in locale (ad esempio Qwen); il ricorso a soluzioni cloud non è necessario. & \mbox{VE-1.8}, \mbox{LG-2.1} \\ \hline
    \mbox{R7-V-O} & Obbligatorio & L'orchestrazione dei modelli linguistici deve avvalersi del framework LangChain4j. & \mbox{Capitolato}, \mbox{VE-1.2}, \mbox{LG-2.1} \\ \hline
    \mbox{R8-V-O} & Obbligatorio & Il motore deve basarsi su una pipeline modulare event-driven con scheduler real-time multi-canale. & \mbox{Capitolato}, \mbox{LG-2.1} \\ \hline
    \mbox{R9-V-O} & Obbligatorio & Il template JSON deve essere definito in funzione delle API fornite per i software aziendali, come estensione delle stesse, con possibilità di introdurre nuove informazioni se necessario. & \mbox{VE-1.3}, \mbox{LG-2.3} \\ \hline
    \mbox{R10-V-O} & Obbligatorio & L'integrazione con i software aziendali (B2B, ERP, CRM, ticketing) deve essere realizzata e verificata sui mock e sulle API dedicate fornite dall'azienda proponente. & \mbox{VE-1.1}, \mbox{LG-2.3} \\ \hline
    \mbox{R11-V-D} & Desiderabile & Per la trascrizione audio e l'OCR dovrebbero essere utilizzati gli strumenti consigliati dall'azienda proponente (Whisper, Google Speech, Tesseract, EasyOCR). & \mbox{Capitolato}, \mbox{LG-1.1} \\ \hline
    \end{longtable}
}